% !TeX root = ../main.tex

\begin{acknowledgements}
时光荏苒，本科阶段的学习即将结束，本论文也在反复修改与完善中接近尾声。在论文完成之际，我想向所有在本科阶段给予我帮助、指导和支持的老师、同学与家人致以诚挚的感谢。

首先，我要衷心感谢我的指导教师谢传龙副教授。大二上学期，在学习谢老师开设的《统计学习》课程之后，我第一次真切地感受到统计学在人工智能时代的重要价值。随后，在学习《深度学习》课程的过程中，我又进一步对深度学习领域产生了浓厚兴趣。毕业论文撰写期间，谢老师在论文选题、研究思路、方法设计、实验推进以及论文修改等方面都给予了我耐心而细致的指导。老师严谨认真的治学态度使我受益匪浅，也让我更加深刻地体会到学术研究需要踏实、细致与规范。

同时，我要感谢文理学院统计系各位老师在本科阶段对我的培养。作为统计系第一届本科生，我们始终感受到老师们对班级同学的关心与重视。无论是在课堂上对知识点的细致讲解，还是在课后对问题的耐心解答，老师们都给予了我们很多帮助。在一门门课程的学习过程中，我逐渐体会到钻研问题本身的乐趣，也慢慢发现，自己不再只是为了应对考试而学习，而是真正对统计学产生了兴趣。

此外，我还要感谢大学四年中结识的同学和朋友们。无论是在图书馆里一起学习、讨论问题，还是在健身房里彼此督促、相互鼓励，抑或是在紧张的学习之余一起打几局游戏放松心情，这些真实而具体的陪伴都构成了我大学生活中十分珍贵的记忆。正是因为有你们一路同行，我的大学生活才更加丰富而充实，也让我在学习、生活和心态上都获得了成长与进步。

最后，我要感谢我的家人。感谢你们一直以来对我学习和生活的理解、支持与陪伴。进入大学以后，面对陌生的环境和全新的学习节奏，我在适应过程中经历过压力、迷茫与不安，也体会过成长、收获与喜悦。课程内容的难度、与家乡不同的饮食习惯以及南方潮湿多雨的气候，都曾让我一时感到无所适从；而在学习逐渐步入正轨之后，拿到奖学金时的开心、日常生活中的点滴进步，以及大学生活里那些平凡却有趣的小事，也同样成为我想第一时间与你们分享的内容。无论是在我倾诉压力和烦恼的时候，还是在我分享收获与喜悦的时候，你们始终都耐心倾听，给予我理解、安慰与鼓励。正是在你们始终如一的陪伴和支持下，我逐渐适应了大学生活，也学会了以更加平和而坚定的心态面对成长中的困难与挑战。

谨以此文，向所有关心、帮助和支持我的人表示衷心感谢。
\end{acknowledgements}
